\honest{Helpless Individuals}{Eliezer Yudkowsky}{2009}

Our grouping instincts are built for 50-person hunter-gatherer bands, so it seems a miracle that large institutions survive at all. Some do: governments with militaries and police that can tax, a form dating from agriculture and extractable surpluses; corporations whose money is controlled from the center, dating from large-scale trade and specialization; and religions, memes grown more virulent in large populations, with threats of damnation, promises of heaven and professional priests. But most large institutions that would serve us, even vital ones, ``simply fail to exist in the first place.'' \nb{Mancur Olson made a similar argument in \textsc{The Logic of Collective Action} (1965): without coercion or ``some other special device,'' large groups do not act in their common interest. Olson's reason was free riding by self-interested people. The post's reason, which follows, is different.}

I realized this from how science is funded. I have found firsthand that it is ``amazingly difficult'' to raise money for science from individuals, unless it concerns a disease with gruesome victims, ``and maybe not even then.'' \nb{The March of Dimes funded Salk's polio vaccine with small gifts, so that exception can be large.} People are prosocial; they give to puppy pounds, and they know science is a great social interest. But a particular project, say on the genetics of trypanotolerance in cattle, is a poor emotional fit for personal giving. It takes years, its press releases are dull, only specialists can do the work, and a picture of the scientist whose salary you pay lacks the impact of the wide-eyed puppy you helped rehome; donors get no sense of immediate accomplishment. I first saw this when I wondered why no individuals pay to test Seth Roberts's obesity hypothesis or Bussard's polywell fusion device, both ``obviously ridiculously underfunded.'' \nb{The polywell work was then funded by the US Navy.}

``It seems to me'' that very few things fit individual giving, and that this is key to why so many individual interests are poorly protected, and why 200 million adult Americans have such trouble supervising the 535 members of Congress.

So science is funded by governments spending money that is not quite their own, by large corporations funding blue-sky research, by grassroots organizations built around affective death spirals that fund science suiting their ideals, and by foundations, which fund science that sounds charitable, like orchestras or modern art, with money ``block-allocated by wealthy individuals to their reputations.'' I give no evidence for that motive. Scientists then fight over the pre-allocated money before grant committees. A project rarely bids directly for society's resources, and the exceptions, like the moon shot or the Manhattan Project, were driven by politicians or the military. I grant that if the public funded science directly, much would be wasted on ``quantum gibberish.''

My conclusion is that science survives as a ``parasite'' on the few large organizations that can exist. Many other projects simply never exist. There are only big taxers, big traders, supermemes, occasional individuals of great power, and a few parasites like science, and coordination on common interests fails because it is ``too non-ancestral a problem when you scale to more than 50 people.''
